\newif\ifglobal

%%%%%%%%%%%%%%%%%%%%%%%
%%% Compile to view %%%
%%%%%%%%%%%%%%%%%%%%%%%

%%%%%%%%%%%%%%%%%%%%%%%%%%%%%%%%%%%%%%%%%%%%%%%%%%%
%%% Readme:                                     %%%
%%% https://www.overleaf.com/read/bwdjvfjksvjy/ %%%
%%%%%%%%%%%%%%%%%%%%%%%%%%%%%%%%%%%%%%%%%%%%%%%%%%%

% >>> M?: marks a module setting number or important notice

% >>> M4: Set {./relative-path} to external files for this module <<<
\def \modulefiles {./55-SDHOST}
% To add an external file, use:
% (Replace '---' with actual filename)
% '\input{\modulefiles/tables/---.tex}' for tables
% '\includegraphics{\modulefiles/figures/---}' for figures
% '\subfile{\modulefiles/---}' for register files


% >>> M1: This module is in global project? <<<
% 'YES' - uncomment; 'NO' - comment out
\globaltrue

\ifglobal
    \documentclass[main\_\_EN.tex]{subfiles}
   
\else
    \documentclass[openany, 10pt]{book}

    % >>> M2: In ./chip-spec-settings/preamble-trm-module, also find
    %         settings for visibility of todo notes, watermark <<<
    \usepackage{preamble-trm-module}

    % >>> M3: Temporary labels in English,
    %         you can add your own labels there <<<
    \myexternaldocument{temp-labels-en}

\fi %\ifglobal


\setcounter{secnumdepth}{4}


% >>> M5: If this module needs any updates in
% './chip-spec-settings/preamble-trm-module.sty'
% (include additional package, etc.), contact your module PM
% or create the file './chip-spec-settings/changelog.tex' make the
% necessary updates yourself and document them in this file <<<

\raggedright


\begin{document}


% Sets language into which pre-defined bits of text are translated
\selectlanguage{english}

% Do not delete this block
\ifglobal
% Add the following front matter if
% this module is outside of global project
\else
    \InputIfFileExists{readme}{}{}
    {\let\clearpage\relax \listoftodos}
    {\let\clearpage\relax \listofdonetodos}
    \tableofcontents
    %\listoftables
    %\listoffigures
\fi


%%%%%%%%%%%%%%%%%%%%%%%%%
%%% TRM Module Begins %%%
%%%%%%%%%%%%%%%%%%%%%%%%%

\hypertarget{sdhost}{}
\chapter{SD/MMC Host Controller (SDHOST)}
\label{mod:sdhost}


\section{Overview}

The \chipname{} memory card interface controller provides a hardware interface between the Advanced Peripheral Bus (APB) and an external memory device. The memory card interface allows the \chipname{} to be connected to SDIO (secure digital I/O) memory cards, MMC (multimedia cards) and devices with a CE-ATA (Consumer Electronics Advanced Transport Architecture) interface. It supports two external cards (Card0 and Card1).
All SD/MMC module interface signals only connect to GPIO pins via GPIO matrix.

\section{Features}
This module supports the following features:

\begin{itemize}
    \item Two external cards
    \item SD memory Card specification V3.0 and V3.01, with a maximum transfer rate of DDR50
    \item MMC: V4.41, V4.5, and V4.51, with a maximum transfer rate of DDR50
    \item CE-ATA: V1.1
    \item 1-bit, 4-bit, and 8-bit modes
\end{itemize}

The SD/MMC controller topology is shown in Figure \ref{fig:sdhost-card-system-topo}. The controller supports two peripherals, but they cannot function at the same time. 

\begin{figure}[H]
    \centering
   \includegraphics[width=0.55\textwidth]{\modulefiles/figures/55-SDHOST-card-system-topo.eps}
    \caption{SD/MMC Controller Topology}
    \label{fig:sdhost-card-system-topo}
\end{figure}

\section{SD/MMC External Interface Signals}

The primary external interface signals, which enable the SD/MMC controller to communicate with an external device, are clock (sdhost\_cclk\_out\_1.eg:card1), command (sdhost\_ccmd\_out\_1), and data signals \\(sdhost\_cdata\_in\_1[7:0]/sdhost\_cdata\_out\_1[7:0]). Additional signals include the card interrupt, card detect, and write protect signals. The direction of each signal is shown in Figure \ref{fig:sdhost-external-interface-signals}. The direction and description of each pin are listed in Table \ref{tab:sdhost-signal-description}.

\begin{figure}[H]
    \centering
    \includegraphics[width=1\textwidth]{\modulefiles/figures/55-SDHOST-external-signals.eps}
    \caption{SD/MMC Controller External Interface Signals}
    \label{fig:sdhost-external-interface-signals}
\end{figure}
	
\begin{longtable}[c]{ |p{5cm}|p{3cm}|p{8cm}| }
\caption{SD/MMC Signal Description} \label{tab:sdhost-signal-description}\\
\hline
\rowcolor{lightgray}
Signal & Direction & Description \\
\hline
\endfirsthead
\hline
\rowcolor{lightgray}
Signal & Direction & Description  \\
\hline
\endhead
\endfoot
sdhost\_cclk\_out & Output & Clock signals for slave device	\\\hline
sdhost\_ccmd & Duplex & Duplex command/response lines  \\\hline
sdhost\_cdata & Duplex & Duplex data read/write lines	\\\hline
sdhost\_card\_detect\_n & Input &  Card detection input line\\\hline
sdhost\_card\_write\_prt & Input & Card write protection status input 	\\\hline
\end{longtable}


\section{Functional Description}
\subsection{SD/MMC Host Controller Architecture}
The SD/MMC host controller consists of two main functional blocks, as shown in Figure \ref{fig:sdhost-block-diagram}:

\begin{itemize}
    \item Bus Interface Unit (BIU): Provides the APB interface for registers, access to RAM data, and DMA data read and write operations.
    \item Card Interface Unit (CIU): Handles external memory card interface protocols and provides clock control.
\end{itemize}

\begin{figure}[H]
    \centering
    \includegraphics[width=0.9\textwidth]{\modulefiles/figures/55-SDHOST-block-diagram}
    \caption{SDIO Host Block Diagram}
    \label{fig:sdhost-block-diagram}
\end{figure}

\subsubsection{Bus Interface Unit (BIU)}
The BIU provides access to registers and RAM data through the Host Interface Unit (HIU). Additionally, it provides a method to access to memory data through a DMA interface. Figure \ref{fig:sdhost-block-diagram} illustrates the internal components of the BIU.

The BIU provides the following functions: 
	
\begin{itemize}
	\item Host interface 
	\item DMA interface 
	\item Interrupt control 
	\item Register access 
	\item FIFO access 
	\item Power-up/pull-up control and card detection
\end{itemize}

\subsubsection{Card Interface Unit (CIU)}
The CIU module implements the dedicated card slot specifications. Within the CIU, the command path control unit and data path control unit are used to interface with the command and data ports, respectively, of the SD/MMC/CE-ATA cards. The CIU also provides clock control. Figure \ref{fig:sdhost-block-diagram} illustrates the internal structure of the CIU, which consists of the following primary functional blocks:
	
\begin{itemize}
	\item Command path
	\item Data path 
	\item SDIO interrupt control 
	\item Clock control 
	\item MUX (multiplexer) and De-MUX (de-multiplexer) unit
\end{itemize}

\subsection{Command Path}
The command path performs the following functions: 

\begin{itemize}
	\item Configures clock parameters 
	\item Configures card command parameters 
	\item Sends commands to card bus (sdhost\_ccmd\_out line) 
	\item Receives responses from card bus (sdhost\_ccmd\_in line) 
	\item Sends responses to BIU 
	\item Drives the P-bit on the command line
\end{itemize}
	
The command path state machine is shown in Figure \ref{fig:sdhost-cmd-state-machine}.


\begin{figure}[H]
    \centering
    \includegraphics[width=0.55\textwidth]{\modulefiles/figures/55-SDHOST-cmd-state-machine.eps}
    \caption{Command Path State Machine}
    \label{fig:sdhost-cmd-state-machine}
\end{figure}

\subsection{Data Path}

The data path block pops RAM data and transmits them on sdhost\_cdata\_out during a write-data transfer, or it receives data on sdhost\_cdata\_in and pushes them into RAM during a read-data transfer. The data path loads new data parameters, i.e., expected data, read/write data transfer, stream/block transfer, block size, byte count, card type, timeout registers, etc., whenever a data transfer command is not in progress. 

If the \hyperref[fielddesc:SDHOSTDATAEXPECTED]{SDHOST\_DATA\_EXPECTED} bit is set in the \hyperref[regdesc:SDHOSTCMDREG]{SDHOST\_CMD\_REG} register, the new command is a data transfer command and the data path starts one of the following operations: 

\begin{itemize}
\item Transmitting data if the \hyperref[fielddesc:SDHOSTREADWRITE]{SDHOST\_READ\_WRITE} bit is 1
\item Receiving data if the \hyperref[fielddesc:SDHOSTREADWRITE]{SDHOST\_READ\_WRITE} bit is 0
\end{itemize}
	
\subsubsection{Data Transmit}

The module starts data transmission two clock cycles after a response for the data write command is received. This occurs even if the command path detects a response error or a cyclic redundancy check (CRC) error in a response. If no response is received from the card until the response timeout, no data will be transmitted. Depending on the value of the \hyperref[fielddesc:SDHOSTTRANSFERMODE]{SDHOST\_TRANSFER\_MODE} bit in the \hyperref[regdesc:SDHOSTCMDREG]{SDHOST\_CMD\_REG} register, the data transmit state machine adds data to the card's data bus in a stream or in block(s).
The data transmit state machine is shown in Figure \ref{fig:sdhost-tx-data-state-machine}.

\begin{figure}[H]
    \centering
    \includegraphics[width=0.7\textwidth]{\modulefiles/figures/55-SDHOST-tx-state-machine.eps}
    \caption{Data Transmit State Machine}
    \label{fig:sdhost-tx-data-state-machine}
\end{figure}

\subsubsection{Data Receive}

The module receives data two clock cycles after the end bit of a data read command, even if the command path detects a response error or a CRC error. If no response is received from the card and a response timeout occurs, the BIU does not receive a signal about the completion of the data transfer. If the command sent by the CIU is an illegal operation for the card, it would prevent the card from starting a read-data transfer, and the BIU will not receive a signal about the completion of the data transfer.

If no data is received by the data timeout, the data path signals a data timeout to the BIU, which marks an end to the data transfer. Based on the value of the \hyperref[fielddesc:SDHOSTTRANSFERMODE]{SDHOST\_TRANSFER\_MODE} bit in the \hyperref[regdesc:SDHOSTCMDREG]{SDHOST\_CMD\_REG} register, the data receive state machine gets data from the card's data bus in a stream or block(s).
The data receive state machine is shown in Figure \ref{fig:sdhost-rx-data-state-machine}.

\begin{figure}[H]
    \centering
    \includegraphics[width=0.7\textwidth]{\modulefiles/figures/55-SDHOST-rx-state-machine.eps}
    \caption{Data Receive State Machine}
    \label{fig:sdhost-rx-data-state-machine}
\end{figure}

\section{Software Restrictions for CIU Operations}
\begin{itemize}
	\item Only one card at a time can be selected to execute a command or data transfer. For example, when data are being transferred to or from a card, a new command must not be issued to another card. A new command, however, can be issued to the same card, allowing it to read the device status or stop the transfer. 
	\item Only one data transfer command can be issued at a time. 
	\item During an open-ended card write operation, if the card clock is stopped because RAM is empty, the software must first fill RAM with data, and start the card clock. Only then can it issue a stop/abort command to the card. 
	\item During an SDIO/Combo card transfer, if the card function is suspended and the software wants to resume the suspended transfer, it must first reset and then release RAM by configuring the \hyperref[fielddesc:SDHOSTFIFORESET]{SDHOST\_FIFO\_RESET} bit, and then issue the resume command as if it were a new data transfer command. 
	\item When issuing card reset commands (CMD0, CMD15 or CMD52\_reset) while a card data transfer is in progress, the software must set the \hyperref[fielddesc:SDHOSTSTOPABORTCMD]{SDHOST\_STOP\_ABORT\_CMD} bit in the \hyperref[regdesc:SDHOSTCMDREG]{SDHOST\_CMD\_REG} register, so that the CIU can stop the data transfer after issuing the card reset command. 
	\item When the controlling bit for end bit error is set in the \hyperref[regdesc:SDHOSTRINTSTSREG]{SDHOST\_RINTSTS\_REG} register, the CIU cannot control SDIO interrupts. In such a case, the software must ignore SDIO interrupts and issue a stop/abort command to the card, so that the card stops transmitting data. 
    \item If the card clock is stopped because RAM is full during a card read, the software must read at least two RAM addresses to restart the card clock. 
	\item Only one CE-ATA device at a time can be selected for a command or data transfer. For example, when data are transferred from a CE-ATA device, a new command must not be sent to another CE-ATA device.
	\item If the CE-ATA device interrupts are enabled (nIEN=0), a new SDHOST\_RW\_BLK command must not be sent to the same device if the execution of a SDHOST\_RW\_BLK command is already in progress. Only the Command Completion Signal Disable (CCSD) command can be sent while waiting for the Command Completion Signal (CCS). 
	\item If, however, the CE-ATA device interrupts are disabled (nIEN=1), a new command can be issued to the same device, allowing it to read status information. 
	\item Open-ended transfers are not supported for CE-ATA devices.
	\item The sdhost\_send\_auto\_stop signal is not supported (software must not set the \hyperref[fielddesc:SDHOSTSENDAUTOSTOP]{SDHOST\_SEND\_AUTO\_STOP} bit) for CE-ATA transfers.
\end{itemize}

After configuring the command start bit to 1, the values of the following registers cannot be changed before a command has been issued: 

\begin{itemize}
	\item CMD - command
	\item CMDARG -  command argument
	\item BYTCNT -  byte count
	\item BLKSIZ -  block size
	\item CLKDIV -  clock divider
	\item CKLENA -  clock enable
	\item CLKSRC -  clock source
	\item TMOUT -   timeout
	\item CTYPE -   card type
\end{itemize}

\section{RAM for Receiving and Transmitting Data}

The submodule RAM is a buffer area for transmitting and receiving data. It can be divided into two units: one is for transmitting data, and the other is for receiving data. The process of transmitting and receiving data can be achieved by the CPU and DMA for reading and writing. The latter method is described in detail in Section \ref{sec:sdhost-linked-list}.

\subsection{TX RAM Module}

There are two ways to enable a write operation: CPU and DMA write.

For CPU write operations  to TX RAM, write the data directly into the \hyperref[fielddesc:SDHOSTBUFFIFOREG]{\hyperref[regdesc:SDHOSTBUFFIFOREG]{SDHOST\_BUFFIFO\_REG}} register via the APB interface. The corresponding FIFO width is 32 bits and the depth is 512.

For DMA write operations to TX RAM,  set up DMA linked list descriptors as described in Section \ref{sec:sdhost-linked-list}.

\subsection{RX RAM Module}

There are two ways to enable a read operation: CPU and DMA read.

For CPU read operations from RX RAM, read the data directly from the \hyperref[fielddesc:SDHOSTBUFFIFOREG]{\hyperref[regdesc:SDHOSTBUFFIFOREG]{SDHOST\_BUFFIFO\_REG}} register via the APB interface. The corresponding FIFO width is 32 bits and the depth is 512.

For DMA read operations from RX RAM,  set up DMA linked list descriptors as described in Section \ref{sec:sdhost-linked-list}.

\section{DMA Linked List}
Each unit in the linked list structure consists of two parts: a descriptor and a data buffer. In other words, each descriptor points to a unique data buffer and to the next descriptor in the linked list. Figure \ref{fig:sdhost-chain-descriptor} shows the linked list.

\begin{figure}[H]
    \centering
    \includegraphics[width=0.7\textwidth]{\modulefiles/figures/55-SDHOST-chain-descriptor.eps}
    \caption{Linked List Structure}
    \label{fig:sdhost-chain-descriptor}
\end{figure}

\section{DMA Descriptor Format}
\label{sec:sdhost-linked-list}
Each descriptor consists of four words as shown in \ref{fig:sdhost-linked-list-structure}. Table \ref{tab:sdhost-DES0}, Table \ref{tab:sdhost-DES1}, Table \ref{tab:sdhost-DES2}, Table \ref{tab:sdhost-DES3} provide the descriptions of these four words.

\begin{figure}[H]
    \centering
    \includegraphics[width=0.7\textwidth]{\modulefiles/figures/55-SDHOST-linkded-list}
    \caption{Descriptor Format}
    \label{fig:sdhost-linked-list-structure}
\end{figure}

The DES0 field contains control and status information.

\begin{longtable}[c]{ |p{3cm}|p{4cm}|p{8cm}| }
\caption{DES0 Descriptor Field} \label{tab:sdhost-DES0}\\
\hline
\rowcolor{lightgray}
Bits &Name & Description \\
\hline
\endfirsthead
\hline
\rowcolor{lightgray}
Bits &Name & Description  \\
\hline
\endhead
\endfoot
\multirow{4}{*}{31} & \multirow{4}{*}{OWNER} & When set to 1, this bit indicates that the descriptor is owned by the DMA Controller. When reset to 0, it indicates that the descriptor is owned by the host. The DMA clears this bit when it completes the data transfer.\\\hline
\clearpage
\multirow{11}{*}{30} & \multirow{11}{*}{CES (Card Error Summary)} & This bit indicates the status of the card's read/write operations. \\
 & & This bit is the logical OR of the following bits in the \hyperref[regdesc:SDHOSTRINTSTSREG]{SDHOST\_RINTSTS\_REG} register.
 \begin{itemize}
     \item{EBE: End Bit Error}
     \item{RTO: Response Timeout}
     \item{RCRC: Response CRC}
     \item{SBE: Start Bit Error}
     \item{DRTO: Data Read Timeout}
     \item{DCRC: Data CRC for Receive}
     \item{RE: Response Error}
 \end{itemize}\vspace{-1.5em}
 \\\hline
29:6 & Reserved & Reserved	\\\hline
\multirow{4}{*}{5} &\multirow{4}{*}{ER (End of Ring)} & When set to 1, this bit indicates that the linked list has reached its final descriptor. The DMA Controller then returns to the base address of the linked list, creating a linked list ring. \\\hline
\multirow{3}{*}{4} &\multirow{3}{*}{\makecell[l]{CH\\ (Second Address Chained)}} & When set to 1, this bit indicates that the second address in the descriptor is the next descriptor address. When this bit is set to 1, BS2 (DES1[25:13]) must be all zeros. \\\hline
\multirow{4}{*}{3} & \multirow{4}{*}{FD (First Descriptor)}& When set to 1, this bit indicates that this descriptor contains the first buffer of the data. If the size of the first buffer is 0, the next descriptor contains the beginning of the data. \\\hline
\multirow{7}{*}{2} & \multirow{7}{*}{LD (Last Descriptor)}&  This bit is associated with the last data block of a DMA transfer. When set to 1, the bit indicates that the buffers pointed by this descriptor are the last buffers of the data. After this descriptor is completed, the remaining byte count is 0. In other words, after the descriptor with the LD bit set is completed, the remaining byte count should be 0.\\\hline
\multirow{3}{*}{1}&\multirow{3}{*}{\makecell[l]{DIC (Disable Interrupt\\ on Completion)}}& When set to 1, this bit prevents the setting of the TI/RI bit of the DMA Status Register (\hyperref[regdesc:SDHOSTIDSTSREG]{SDHOST\_IDSTS\_REG}) for the data that ends in the buffer pointed by this descriptor. \\\hline
0 & Reserved& Reserved \\\hline
\end{longtable}


The DES1 field contains the buffer size.


\begin{longtable}[c]{ |p{3cm}|p{4cm}|p{8cm}| }
\caption{DES1 Descriptor Field} \label{tab:sdhost-DES1}\\
\hline
\rowcolor{lightgray}
Bits &Name & Description \\
\hline
\endfirsthead
\hline
\rowcolor{lightgray}
Bits &Name & Description  \\
\hline
\endhead
\endfoot

31:26 & Reserved& Reserved  \\\hline
25:13 & Reserved& Reserved  \\\hline
\multirow{4}{*}{12:0}  & \multirow{4}{*}{BS (Buffer Size)} &  These bits indicate the size of the data buffer byte size, which must be a multiple of four. When the buffer size is not a multiple of four, the resulting behavior is undefined. This field must not be zero.    \\\hline

\end{longtable}


The DES2 field contains the address pointer to the data buffer.\\

\begin{longtable}[c]{ |p{3cm}|p{4cm}|p{8cm}| }
\caption{DES2 Descriptor Field} \label{tab:sdhost-DES2}\\
\hline
\rowcolor{lightgray}
Bits &Name & Description \\
\hline
\endfirsthead
\hline
\rowcolor{lightgray}
Bits &Name & Description  \\
\hline
\endhead
\endfoot
\multirow{2}{*}{31:0}& \multirow{2}{*}{Buffer Address Pointer} & These bits indicate the address of the data buffer. The buffer address must be word-aligned.   \\\hline

\end{longtable}


The DES3 field contains the address pointer to the next descriptor if the present descriptor is not the last one in a linked list structure.

\begin{longtable}[c]{ |p{3cm}|p{4cm}|p{8cm}| }
\caption{DES3 Descriptor Field} \label{tab:sdhost-DES3}\\
\hline
\rowcolor{lightgray}
Bits & Name & Description \\
\hline
\endfirsthead
\hline
\rowcolor{lightgray}
Bits & Name & Description  \\
\hline
\endhead
\endfoot
\multirow{4}{*}{31:0}& \multirow{4}{*}{\makecell[l]{Next Descriptor Address}}& If CH (DES0[4]) is set to 1, these bits contains the pointer to the next descriptor. \\
&&If this is not the last descriptor in a linked list structure, bits 1 and 0 must be zero, i.e., DES3[1:0] = 0.\\\hline
\end{longtable}


\section{Programming Procedures}

\subsection{Initializing Registers}

To initialize registers, perform the following steps:

\begin{enumerate}
\item Write to the control register \hyperref[regdesc:SDHOSTCTRLREG]{SDHOST\_CTRL\_REG}, the raw interrupt register \hyperref[regdesc:SDHOSTRINTSTSREG]{SDHOST\_RINTSTS\_REG}, and SDIO interrupt mask register \hyperref[regdesc:SDHOSTINTMASKREG]{SDHOST\_INTMASK\_REG} to clear pending interrupts and configure controlling and interrupt parameters.
\item Write clock divider configuration register \hyperref[regdesc:SDHOSTCLKDIVREG]{SDHOST\_CLKDIV\_REG}, clock source selection register \hyperref[fielddesc:SDHOSTCLKSRCREG]{\hyperref[regdesc:SDHOSTCLKSRCREG]{SDHOST\_CLKSRC\_REG}}, and clock enable register \hyperref[regdesc:SDHOSTCLKENAREG]{SDHOST\_CLKENA\_REG} to configure the card clock.
\item Write other configuration registers based on card parameters. For example, configure 
    \begin{itemize}
        \item card bus width (\hyperref[regdesc:SDHOSTCTYPEREG]{SDHOST\_CTYPE\_REG})
        \item user ID (\hyperref[fielddesc:SDHOSTUSRIDREG]{\hyperref[regdesc:SDHOSTUSRIDREG]{SDHOST\_USRID\_REG}})
        \item
        UHS mode voltage and DDR (\hyperref[regdesc:SDHOSTUHSREG]{SDHOST\_UHS\_REG})
        \item
        EMMC mode start bit (\hyperref[regdesc:SDHOSTEMMCDDRREG]{SDHOST\_EMMCDDR\_REG})
        \item SDIO mode (\hyperref[regdesc:SDHOSTCLKEDGESELREG]{SDHOST\_CLK\_EDGE\_SEL\_REG})
        \item timeout value (\hyperref[regdesc:SDHOSTTMOUTREG]{SDHOST\_TMOUT\_REG})
        \item debounce filter (\hyperref[regdesc:SDHOSTDEBNCEREG]{SDHOST\_DEBNCE\_REG})
        \item internal phase shift (\hyperref[regdesc:SDHOSTENSHIFTREG]{SDHOST\_ENSHIFT\_REG})
    \end{itemize}
     You can reset the card with \hyperref[regdesc:SDHOSTRSTNREG]{SDHOST\_RST\_N\_REG}.
\end{enumerate}

\subsection{Sending Commands}

To send commands, perform the following steps:

\begin{enumerate}
\item Write the command argument register \hyperref[fielddesc:SDHOSTCMDARGREG]{\hyperref[regdesc:SDHOSTCMDARGREG]{SDHOST\_CMDARG\_REG}} with the appropriate command argument parameter.
\item For data transfer, write the data size register \hyperref[regdesc:SDHOSTBYTCNTREG]{SDHOST\_BYTCNT\_REG}, the block size register \hyperref[regdesc:SDHOSTBLKSIZREG]{SDHOST\_BLKSIZ\_REG} register, and the card threshold control register \hyperref[regdesc:SDHOSTCARDTHRCTLREG]{SDHOST\_CARDTHRCTL\_REG} to configure data transfer parameters.
\item Data transfer can be achieved by DMA or FIFO, and DMA is preferable. 
\begin{itemize}
    \item To transmit and receive data via DMA, configure the DMA Controller according to Section \ref{sec:sdhost-dma-init}, Section \ref{sec:sdhost-dma-tx}, and Section \ref{sec:sdhost-dma-rx}.
    \item To transmit and receive data via FIFO, configure the FIFO threshold via FIFO configuration register \hyperref[regdesc:SDHOSTFIFOTHREG]{SDHOST\_FIFOTH\_REG}. For writing to a card, software should write data to the FIFO via \hyperref[fielddesc:SDHOSTBUFFIFOREG]{\hyperref[regdesc:SDHOSTBUFFIFOREG]{SDHOST\_BUFFIFO\_REG}} continuously during command execution.
\end{itemize}
\item Configure and start the command via the \hyperref[regdesc:SDHOSTCMDREG]{SDHOST\_CMD\_REG}.
\item Wait until the \hyperref[fielddesc:SDHOSTSTARTCMD]{SDHOST\_START\_CMD} bit in the \hyperref[regdesc:SDHOSTCMDREG]{SDHOST\_CMD\_REG} register is cleared by hardware upon command reception. When hardware cannot receive the command, it will set the 12th bit HLE in the \hyperref[regdesc:SDHOSTRINTSTSREG]{SDHOST\_RINTSTS\_REG} register to 1.
\item Read the \hyperref[regdesc:SDHOSTRINTSTSREG]{SDHOST\_RINTSTS\_REG} register or enable interrupts to check command execution status. If needed, read the  \hyperref[regdesc:SDHOSTRESP0REG]{SDHOST\_RESP\regindex{n}\_REG} (\regindex{n} ranges from 0 \verb+~+ 3) register to get command response from the card.
\item If data is read from the card via FIFO, software can get the data returned by the card via reading the \hyperref[fielddesc:SDHOSTBUFFIFOREG]{\hyperref[regdesc:SDHOSTBUFFIFOREG]{SDHOST\_BUFFIFO\_REG}} register.
\end{enumerate}

Please refer to the card specification for the command format, and configure register field values depending on the command format.

\subsection{Initializing DMA}\label{sec:sdhost-dma-init}

To initialize the DMA Controller, perform the following steps:

\begin{enumerate}
\item  Write to the DMA bus mode register (\hyperref[regdesc:SDHOSTBMODREG]{SDHOST\_BMOD\_REG}) to configure the host bus's access parameters. 
\item Write to the DMA interrupt enable register (\hyperref[regdesc:SDHOSTIDINTENREG]{SDHOST\_IDINTEN\_REG}) to mask unnecessary interrupt causes. 
\item Create either a transmit or receive linked list. Then write the starting address of the first descriptor in the linked list to the DMA Controller's linked list base address register (\hyperref[fielddesc:SDHOSTDBADDRREG]{\hyperref[regdesc:SDHOSTDBADDRREG]{SDHOST\_DBADDR\_REG}}). 
\item The DMA Controller engine proceeds to get descriptors from the linked list.
\end{enumerate}

\subsection{Initializing DMA Transmission}\label{sec:sdhost-dma-tx}

To initialize DMA transmission, perform the following steps:

\begin{enumerate}
\item The host sets up the descriptor fields (DES0-DES3) for transmission, and sets the OWNER bit (DES0[31]) to 1. The host also prepares the data buffer. 
\item The host writes the write data command to the command register \hyperref[regdesc:SDHOSTCMDREG]{SDHOST\_CMD\_REG}. 
\item The host sets the required transmit threshold via the \hyperref[fielddesc:SDHOSTTXWMARK]{SDHOST\_TX\_WMARK} field in the \hyperref[regdesc:SDHOSTFIFOTHREG]{SDHOST\_FIFOTH\_REG} register. 
\item The DMA Controller engine fetches the descriptor and checks the OWNER bit. If the OWNER bit is not set, it means that the host owns the descriptor. In this case, the DMA Controller enters a suspend state and asserts the Descriptor Unable interrupt via the \hyperref[regdesc:SDHOSTIDSTSREG]{SDHOST\_IDSTS\_REG} register. The host then needs to release the DMA Controller by writing any value to \hyperref[regdesc:SDHOSTPLDMNDREG]{SDHOST\_PLDMND\_REG}. 
\item The DMA Controller waits for the Command Done (CMDD) bit in the \hyperref[regdesc:SDHOSTRINTSTSREG]{SDHOST\_RINTSTS\_REG} register to be set to 1 with no errors from BIU, which indicates that a transfer has completed. 
\item The DMA Controller engine waits for a DMA interface request from BIU. This request is generated based on the configured transmit threshold value. For the last byte of data that cannot be accessed using a burst, a single transfer is performed on the AHB bus. 
\item The DMA Controller fetches the transmit data from the data buffer in the host memory and transfers them to RAM in preparation for transmission to the card. 
\item When data span across multiple descriptors, the DMA Controller fetches the next descriptor and continues with its operation with the next descriptor. The Last Descriptor bit indicates whether the data span multiple descriptors or not.
\item When data transmission is complete, the status information is updated in the \hyperref[regdesc:SDHOSTIDSTSREG]{SDHOST\_IDSTS\_REG} register by setting the \hyperref[fielddesc:SDHOSTIDSTSTI]{SDHOST\_IDSTS\_TI} bit to 1 if it has already been enabled. Also, the OWNER bit is cleared by the DMA Controller by updating the DES0 field. 
\end{enumerate}

\subsection{Initializing DMA Reception}\label{sec:sdhost-dma-rx}

To initialize DMA reception, perform the following steps:

\begin{enumerate}
\item The host sets up the descriptor fields (DES0-DES3) for reception, and sets the OWNER bit (DES0[31]) to 1.
\item The host writes the read data command in the \hyperref[regdesc:SDHOSTCMDREG]{SDHOST\_CMD\_REG} register in BIU. 
\item The host sets the required receive threshold level via the \hyperref[fielddesc:SDHOSTRXWMARK]{SDHOST\_RX\_WMARK} field in the \hyperref[regdesc:SDHOSTFIFOTHREG]{SDHOST\_FIFOTH\_REG} register. 
\item The DMA Controller engine fetches the descriptor and checks the OWNER bit. If the OWNER bit is not set, it means that the host owns the descriptor. In this case, the DMA enters suspend state and asserts the Descriptor Unable interrupt via the \hyperref[regdesc:SDHOSTIDSTSREG]{SDHOST\_IDSTS\_REG} register. The host then needs to release the DMA Controller by writing any value to \hyperref[regdesc:SDHOSTPLDMNDREG]{SDHOST\_PLDMND\_REG}. 
\item The DMA Controller waits for the Command Done (CMDD) bit in the \hyperref[regdesc:SDHOSTRINTSTSREG]{SDHOST\_RINTSTS\_REG} register to be set to 1 with no errors from BIU, which indicates that a reception has been completed. 
\item The DMA Controller engine waits for a DMA interface request from BIU. This request is generated based on the configured receive threshold value. For the last byte of data that cannot be accessed using a burst, a single transfer is performed on the AHB bus. 
\item The DMA Controller fetches the data from RAM and transfers them to the host memory. 
\item When data span across multiple descriptors, the DMA Controller fetches the next descriptor and continues with its operation with the next descriptor. The Last Descriptor bit indicates whether the data span multiple descriptors or not. 
\item When data reception is complete, the status information is updated in the \hyperref[regdesc:SDHOSTIDSTSREG]{SDHOST\_IDSTS\_REG} register by setting the \hyperref[fielddesc:SDHOSTIDSTSRI]{SDHOST\_IDSTS\_RI} bit to 1, if it has already been enabled. Also, the OWNER bit is cleared by the DMA Controller by updating the DES0 field.
\end{enumerate}

\section{Clock Phase Selection}\label{sec:sdhost-phase-selection}

The phase of the SD/MMC Host Controller clock source is configurable, so as to adjust the setup time sequence for the input or output data signals.

The clock source of the SD/MMC Host Controller can be a high-performance clock at a high frequency or a low-power clock at a low frequency. When the clock source is a high-performance clock, the internal signal clock phase, output signal driving clock phase, and input signal sampling clock phase can be configured via the \hyperref[regdesc:SDHOSTDLLCLKCONFREG]{SDHOST\_DLL\_CLK\_CONF\_REG}, in the unit of 1/64 clock source period. When the clock source is a low-power clock, the internal signal clock phase, output signal driving clock phase, and input signal sampling clock phase can be configured via the system clock register. The phase options are 0 degrees, 90 degrees, 180 degrees, and 270 degrees, or in other words the unit is 1/4 clock source period.

For details, please see Chapter \ref{mod:resclk} \textit{\nameref{mod:resclk}}.

%%%%%%%%%%%%%%%%%%%%%
%%% Interrupts %%%
%%%%%%%%%%%%%%%%%%%%%

\section{Interrupt}

The SD/MMC Host Controller can generate interrupt signal SDIO\_HOST\_INTR triggered by the following interrupt sources, and send the interrupt signals to the \hyperref[mod:intmtrx]{Interrupt Matrix}.

\begin{itemize}
\item SDIO\_INT: Interrupts from SDIO cards
\item EBE: Triggered when an error in end bit occurs, or no data CRC received 
\item ACD: Triggered when an automatically sent command has been executed
\item SBE/BCI: Triggered when an error in the start bit occurs
\item HLE: Triggered when a write error occurs during hardware-lock period
\item FRUN: Triggered when the FIFO is empty or full
\item HTO: Triggered when the host does not fill data before timeout period
\item DRTO: Triggered when a data read operation times out
\item RTO: Triggered when card response times out
\item DCRC: Triggered when data CRC error occurs
\item RCRC: Triggered when a card response CRC error occurs
\item RXDR: Triggered during read operations from card when FIFO level is greater than Receive-Threshold level
\item TXDR: Triggered during write operations to card when FIFO level reaches less than or equal to Transmit-Threshold level.
\item DTO: Triggered when data transfer completes
\item CMDD: Triggered when a command has been executed
\item RE: Triggered when response error occurs
\item CD: Triggered when a card is detected
\item IDSTS\_RI: Triggered when data pointed by a DMA descriptor list have been received
\item IDSTS\_TI: Triggered when data pointed by a DMA descriptor list have been transmitted
\item IDSTS\_FBE: Triggered when a fatal bus error occurs and the host is stopped during data transmission or reception. When triggered, the DMA disables all bus access.
\item IDSTS\_DU: Triggered when a DMA descriptor list is not available because the OWNER bit is 0 (DES0 [31] = 0)
\end{itemize}

The \hyperref[regdesc:SDHOSTRINTSTSREG]{SDHOST\_RINTSTS\_REG} and \hyperref[regdesc:SDHOSTIDSTSREG]{SDHOST\_IDSTS\_REG} register contains all the bits that might cause an interrupt. The \hyperref[regdesc:SDHOSTINTMASKREG]{SDHOST\_INTMASK\_REG} and \hyperref[regdesc:SDHOSTIDINTENREG]{SDHOST\_IDINTEN\_REG} register contains an enable bit for each of the events that can cause an interrupt. 

There are three interrupt summary bits in the \hyperref[regdesc:SDHOSTIDSTSREG]{SDHOST\_IDSTS\_REG} register: the card error interrupt summary bit (bit 5 \hyperref[fielddesc:SDHOSTIDSTSCES]{SDHOST\_IDSTS\_CES}), the normal interrupt summary bit (bit 8 \hyperref[fielddesc:SDHOSTIDSTSNIS]{SDHOST\_IDSTS\_NIS}), and the abnormal interrupt summary bit (bit 9 \hyperref[fielddesc:SDHOSTIDSTSAIS]{SDHOST\_IDSTS\_AIS}). Interrupts are cleared by writing 1 to the corresponding bit. When all the enabled interrupts within a group are cleared, the corresponding summary bit is also cleared.

Interrupts are not queued. If another interrupt event occurs before the driver has responded to the previous interrupt, no additional interrupts are generated. For example, the \hyperref[fielddesc:SDHOSTIDSTSRI]{SDHOST\_IDSTS\_RI} indicates that one or more data were transferred to the host buffer. 

An interrupt is generated only once for simultaneous multiple events. The driver must scan the \hyperref[regdesc:SDHOSTRINTSTSREG]{SDHOST\_RINTSTS\_REG} and \hyperref[regdesc:SDHOSTIDSTSREG]{SDHOST\_IDSTS\_REG} register for the interrupt cause.

\begin{tiplisting}
For definitions of \textit{interrupt}, \textit{interrupt signal}, \textit{interrupt source}, and their correlations, please refer to Chapter \ref{mod:intmtrx} \textit{\nameref{mod:intmtrx}} > Section \ref{sec:intmtx-int-terminology} \textit{\nameref{sec:intmtx-int-terminology}}.
\end{tiplisting}

% >>> If your module has no registers, delete sample text
%       for Sections Register Summary and Registers below <<<

%%%%%%%%%%%%%%%%%%%%%%%%
%%% Register Summary %%%
%%%%%%%%%%%%%%%%%%%%%%%%

% >>> Update the sample text below and insert
%     the info on register summary into the subfile <<<

\newpage
\hypertarget{sdhost-reg-summ}{}
\section{Register Summary}
\label{sec:sdhost-reg-summ}

The addresses in this section are relative to the SD/MMC Host Controller base address provided in Table \ref{tab:sysmem-base-address} in Chapter \ref{mod:sysmem} \textit{\nameref{mod:sysmem}}.

The abbreviations given in Column \textbf{Access} are explained in Section \hyperref[glossary-access-types]{\textit{Access Types for Registers}}.

\subfile{\modulefiles/55-SDHOST-reg-summ__EN}


%%%%%%%%%%%%%%%%%
%%% Registers %%%
%%%%%%%%%%%%%%%%%

% >>> Update the sample text below and insert
%      the info on registers into the subfile <<<

\newpage
\section{Registers}

The addresses in this section are relative to the SD/MMC Host Controller base address provided in Table \ref{tab:sysmem-base-address} in Chapter \ref{mod:sysmem} \textit{\nameref{mod:sysmem}}.

For how to program reserved fields, please refer to Section \hyperref[programming-reserved-field]{Programming Reserved Register Field}.

\subfile{\modulefiles/55-SDHOST-reg__EN}


\end{document}
